% ─────────────────────────────────────────────────────────────────────────────
%  main.tex  –  Documento raíz
%  Compilar: pdflatex main.tex  (×2)  +  bibtex main  +  pdflatex main.tex (×2)
% ─────────────────────────────────────────────────────────────────────────────
\documentclass[12pt,a4paper,twoside,openright]{report}

% ─────────────────────────────────────────────────────────────────────────────
%  preamble.tex  –  Paquetes y configuración global
% ─────────────────────────────────────────────────────────────────────────────

% ── Codificación y lengua ────────────────────────────────────────────────────
\usepackage[utf8]{inputenc}
\usepackage[T1]{fontenc}
\usepackage[spanish,es-noshorthands]{babel}

% ── Tipografía ───────────────────────────────────────────────────────────────
\usepackage{lmodern}
\usepackage{microtype}

% ── Matemáticas ──────────────────────────────────────────────────────────────
\usepackage{amsmath,amssymb,amsthm}
\usepackage{mathtools}

% ── Geometría de página ──────────────────────────────────────────────────────
\usepackage[
  a4paper,
  top=2.5cm, bottom=2.5cm,
  left=3cm,  right=2.5cm,
  headheight=15pt
]{geometry}

% ── Gráficos y diagramas ─────────────────────────────────────────────────────
\usepackage{graphicx}
\usepackage{tikz}
\usetikzlibrary{shapes.geometric, arrows.meta, positioning, fit, calc, backgrounds}
\usepackage{pgfplots}
\pgfplotsset{compat=1.18}

% ── Código fuente ────────────────────────────────────────────────────────────
\usepackage{listings}
\usepackage{xcolor}

\definecolor{codegreen}{rgb}{0.13,0.55,0.13}
\definecolor{codegray}{rgb}{0.5,0.5,0.5}
\definecolor{codepurple}{rgb}{0.58,0,0.82}
\definecolor{codeblue}{rgb}{0.0,0.2,0.7}
\definecolor{backcolour}{rgb}{0.97,0.97,0.97}
\definecolor{keywordcolor}{rgb}{0.0,0.2,0.7}

\lstdefinestyle{java}{
  backgroundcolor=\color{backcolour},
  commentstyle=\color{codegreen},
  keywordstyle=\color{keywordcolor}\bfseries,
  numberstyle=\tiny\color{codegray},
  stringstyle=\color{codepurple},
  basicstyle=\ttfamily\footnotesize,
  breakatwhitespace=false,
  breaklines=true,
  captionpos=b,
  keepspaces=true,
  numbers=left,
  numbersep=5pt,
  showspaces=false,
  showstringspaces=false,
  showtabs=false,
  tabsize=4,
  language=Java,
  morekeywords={@pre,@post,@invariant,@variant,result},
  frame=single,
  rulecolor=\color{gray!40},
}

\lstdefinestyle{asm}{
  backgroundcolor=\color{backcolour},
  commentstyle=\color{codegreen}\itshape,
  keywordstyle=\color{codeblue}\bfseries,
  basicstyle=\ttfamily\footnotesize,
  breaklines=true,
  captionpos=b,
  numbers=left,
  numberstyle=\tiny\color{codegray},
  numbersep=5pt,
  frame=single,
  rulecolor=\color{gray!40},
  language={[x86masm]Assembler},
}

\lstdefinestyle{smt}{
  backgroundcolor=\color{backcolour},
  commentstyle=\color{codegreen}\itshape,
  keywordstyle=\color{codepurple}\bfseries,
  basicstyle=\ttfamily\footnotesize,
  breaklines=true,
  captionpos=b,
  numbers=left,
  numberstyle=\tiny\color{codegray},
  numbersep=5pt,
  frame=single,
  rulecolor=\color{gray!40},
  morekeywords={declare-fun,assert,check-sat,get-model,push,pop,exit,set-logic,not,and,or,=>,>=,<=,+,-,*,/},
}

\lstdefinestyle{cpp}{
  backgroundcolor=\color{backcolour},
  commentstyle=\color{codegreen}\itshape,
  keywordstyle=\color{keywordcolor}\bfseries,
  basicstyle=\ttfamily\footnotesize,
  breaklines=true,
  captionpos=b,
  numbers=left,
  numberstyle=\tiny\color{codegray},
  numbersep=5pt,
  frame=single,
  rulecolor=\color{gray!40},
  language=C++,
}

\lstset{style=java}   % estilo por defecto

% ── Tablas ───────────────────────────────────────────────────────────────────
\usepackage{booktabs}
\usepackage{tabularx}
\usepackage{multirow}
\usepackage{array}

% ── Hiperenlaces ─────────────────────────────────────────────────────────────
\usepackage[
  colorlinks=true,
  linkcolor=black,
  citecolor=codeblue,
  urlcolor=codeblue,
  pdftitle={Compilador JavaSubset → x86-64 con Verificación Formal},
  pdfauthor={nickhernd},
]{hyperref}

% ── Encabezados y pies de página ─────────────────────────────────────────────
\usepackage{fancyhdr}
\pagestyle{fancy}
\fancyhf{}
\fancyhead[LE,RO]{\thepage}
\fancyhead[RE]{\small\nouppercase{\leftmark}}
\fancyhead[LO]{\small\nouppercase{\rightmark}}
\renewcommand{\headrulewidth}{0.4pt}

% ── Entornos matemáticos ─────────────────────────────────────────────────────
\theoremstyle{definition}
\newtheorem{definicion}{Definición}[chapter]
\newtheorem{ejemplo}{Ejemplo}[chapter]
\newtheorem{regla}{Regla}[chapter]

\theoremstyle{plain}
\newtheorem{teorema}{Teorema}[chapter]
\newtheorem{lema}{Lema}[chapter]

\theoremstyle{remark}
\newtheorem{nota}{Nota}[chapter]
\newtheorem{observacion}{Observación}[chapter]

% ── Pseudocódigo (sin paquete algorithm2e) ───────────────────────────────────
\lstdefinestyle{pseudo}{
  backgroundcolor=\color{backcolour},
  basicstyle=\ttfamily\footnotesize,
  breaklines=true,
  captionpos=b,
  frame=single,
  rulecolor=\color{gray!40},
  numbers=left,
  numberstyle=\tiny\color{codegray},
  numbersep=5pt,
}

% ── Miscelánea ───────────────────────────────────────────────────────────────
\usepackage{enumitem}
\usepackage{caption}
\usepackage{subcaption}
\usepackage{float}
\usepackage{csquotes}
% backnaur no disponible; gramáticas se presentan en lstlisting

% ── Macros del proyecto ──────────────────────────────────────────────────────
\newcommand{\jss}{\texttt{JavaSubset}}
\newcommand{\IR}{\textsc{ir}}
\newcommand{\WP}[2]{\mathrm{WP}(#1,\,#2)}
\newcommand{\SP}[2]{\mathrm{SP}(#1,\,#2)}
\newcommand{\hoare}[3]{\{#1\}\;#2\;\{#3\}}
\newcommand{\subst}[3]{#1[#3/#2]}           % P[e/x]
\newcommand{\true}{\mathtt{true}}
\newcommand{\false}{\mathtt{false}}
\newcommand{\reg}[1]{\texttt{\%#1}}


% ─────────────────────────────────────────────────────────────────────────────
\begin{document}

% ── Portada ──────────────────────────────────────────────────────────────────
\begin{titlepage}
  \centering
  \vspace*{2cm}

  {\LARGE\bfseries Compilador \texttt{JavaSubset} $\rightarrow$ x86-64}\\[0.5em]
  {\large con Verificación Formal mediante Lógica de Hoare y Z3}

  \vspace{2cm}
  \rule{\textwidth}{0.5pt}
  \vspace{0.5cm}

  {\large\itshape Documentación Técnica del Proyecto}\\[0.3em]
  {\normalsize Diseño de Compiladores — Fases 1 a 5}

  \vspace{0.5cm}
  \rule{\textwidth}{0.5pt}

  \vfill

  \begin{tabular}{rl}
    \textbf{Autor:}      & nickhernd \\[0.3em]
    \textbf{Repositorio:} & \url{https://github.com/nickhernd/pl} \\[0.3em]
    \textbf{Lenguaje:}   & C++11 (Flex + Bison) \\[0.3em]
    \textbf{Herramientas:} & GCC, Z3 SMT Solver \\[0.3em]
    \textbf{Fecha:}      & \today \\
  \end{tabular}

  \vspace{2cm}
\end{titlepage}

% ── Resumen ──────────────────────────────────────────────────────────────────
\begin{abstract}
\noindent
Este documento describe el diseño e implementación de un compilador completo
para un subconjunto de Java con destino nativo x86-64. El compilador abarca
cinco milestones: análisis léxico-sintáctico con Flex y Bison, análisis
semántico con tablas de símbolos y verificación de tipos, generación de código
intermedio en formato de cuádruplas de tres direcciones, backend x86-64 con
generación de ensamblador AT\&T y asignación de registros mediante el
algoritmo de coloración de grafos de Chaitin-Briggs, y finalmente un módulo
de verificación formal basado en la Lógica de Hoare y el Cálculo de
Precondiciones Más Débiles (\emph{Weakest Precondition}, WP) con exportación
a formato SMT-Lib2 para el solver Z3.

El compilador es capaz de procesar programas Java anotados con especificaciones
formales (\texttt{@pre}, \texttt{@post}, \texttt{@invariant}, \texttt{@variant})
y generar automáticamente condiciones de verificación cuya validez es comprobada
por Z3, detectando así tanto corrección de programas como errores de
division por cero y no terminación.
\end{abstract}

\tableofcontents
\listoffigures
\listoftables
\lstlistoflistings

% ── Capítulos ─────────────────────────────────────────────────────────────────
% ─────────────────────────────────────────────────────────────────────────────
%  Capítulo 1 – Introducción
% ─────────────────────────────────────────────────────────────────────────────
\chapter{Introducción}
\label{chap:intro}

\section{Motivación y objetivos}

El diseño e implementación de un compilador desde cero constituye uno de los
proyectos más completos en el ámbito de la Ingeniería del Software, pues pone
en juego simultáneamente teoría de lenguajes formales, estructuras de datos
avanzadas, generación de código máquina y, en este caso, lógica matemática
aplicada a la verificación de programas.

El proyecto \textbf{JavaSubset Compiler} tiene como objetivo construir un
compilador completo para un subconjunto estáticamente tipado de Java con destino
a código nativo x86-64. El proceso se organiza en cinco fases o \emph{milestones}:

\begin{enumerate}[label=\textbf{Fase \arabic*.}]
  \item \textbf{Frontend}: análisis léxico y sintáctico, construcción del AST.
  \item \textbf{Análisis semántico}: resolución de nombres, comprobación de tipos.
  \item \textbf{Representación intermedia}: traducción a código de tres direcciones.
  \item \textbf{Backend x86-64}: generación de ensamblador, asignación de registros.
  \item \textbf{Verificación formal}: Lógica de Hoare, cálculo WP, integración con Z3.
\end{enumerate}

\section{Subconjunto del lenguaje}

El lenguaje fuente es un subconjunto tipado de Java con las siguientes
características:

\begin{table}[H]
  \centering
  \caption{Características del subconjunto Java soportado}
  \label{tab:lenguaje}
  \begin{tabularx}{\textwidth}{lX}
    \toprule
    \textbf{Categoría}   & \textbf{Construcciones soportadas} \\
    \midrule
    Tipos primitivos     & \texttt{int}, \texttt{double}, \texttt{boolean} \\
    Declaraciones        & Variables locales, constantes \texttt{final} \\
    Operadores           & $+$, $-$, $*$, $/$ , $\%$, comparación, lógicos \\
    Control de flujo     & \texttt{if}/\texttt{else}, \texttt{while} \\
    Métodos              & Estáticos sin parámetros, \texttt{return} \\
    E/S estándar         & \texttt{System.out.print/println}, \texttt{Scanner} \\
    Anotaciones formales & \texttt{@pre}, \texttt{@post}, \texttt{@invariant}, \texttt{@variant} \\
    \bottomrule
  \end{tabularx}
\end{table}

\section{Herramientas utilizadas}

\begin{description}
  \item[\textbf{Flex 2.6}] Generador de analizadores léxicos a partir de
    expresiones regulares. Produce la función \texttt{yylex()} en C.
  \item[\textbf{GNU Bison 3.x}] Generador de analizadores sintácticos LALR(1).
    Produce el autómata de análisis y las acciones semánticas en C.
  \item[\textbf{GCC / G++ 11}] Compilador del código C/C++ generado.
    Se usa también para ensamblar y enlazar el código x86-64 producido.
  \item[\textbf{Z3 4.x}] Solver SMT de Microsoft Research. Recibe fórmulas
    en formato SMT-Lib2 y determina su satisfacibilidad (SAT/UNSAT).
\end{description}

\section{Estructura del documento}

El documento sigue la estructura del pipeline del compilador:
el capítulo~\ref{chap:arquitectura} presenta la vista global del sistema;
los capítulos~\ref{chap:frontend}–\ref{chap:backend} describen cada fase
de compilación; el capítulo~\ref{chap:optimizacion} detalla las
optimizaciones sobre la RI; el capítulo~\ref{chap:registro} explica la
asignación de registros; el capítulo~\ref{chap:verificacion} desarrolla
la verificación formal; y el capítulo~\ref{chap:pruebas} recoge los
resultados experimentales.

% ─────────────────────────────────────────────────────────────────────────────
%  Capítulo 2 – Arquitectura del compilador
% ─────────────────────────────────────────────────────────────────────────────
\chapter{Arquitectura del compilador}
\label{chap:arquitectura}

\section{Vista global del pipeline}

El compilador sigue la arquitectura clásica en fases consecutivas. Cada fase
transforma una representación del programa en otra de mayor nivel de abstracción
hacia la máquina destino, o bien extrae información (tabla de símbolos, tipos)
que se usa en fases posteriores.

\begin{figure}[H]
\centering
\begin{tikzpicture}[
  node distance=0.6cm and 1.8cm,
  box/.style={rectangle, draw, rounded corners=3pt,
              minimum width=3.8cm, minimum height=0.8cm,
              align=center, font=\small\ttfamily, fill=blue!8},
  arr/.style={-Stealth, thick},
  lbl/.style={font=\scriptsize\itshape, text=gray}
]

\node[box] (src)    {Fuente .java};
\node[box, below=of src]  (lex)    {Analizador Lexico\\(lexer.l)};
\node[box, below=of lex]  (par)    {Analizador Sintactico\\(parser.y)};
\node[box, below=of par]  (sem)    {Analizador Semantico\\(SemanticVisitor)};
\node[box, below=of sem]  (ir)     {Generacion de RI\\(IRBuilder)};
\node[box, below=of ir]   (opt)    {Optimizacion RI\\(IROptimizer)};
\node[box, below=of opt]  (reg)    {Asig. de Registros\\(RegisterAllocator)};
\node[box, below=of reg]  (x86)    {Generacion x86-64\\(X86Generator)};
\node[box, below=of x86]  (asm)    {Ensamblador .s};

\draw[arr] (src) -- (lex)
  node[lbl, right, midway] {codigo fuente};
\draw[arr] (lex) -- (par)
  node[lbl, right, midway] {tokens};
\draw[arr] (par) -- (sem)
  node[lbl, right, midway] {AST};
\draw[arr] (sem) -- (ir)
  node[lbl, right, midway] {AST tipado};
\draw[arr] (ir) -- (opt)
  node[lbl, right, midway] {cuadruplas IR};
\draw[arr] (opt) -- (reg)
  node[lbl, right, midway] {IR optimizado};
\draw[arr] (reg) -- (x86)
  node[lbl, right, midway] {RegMap};
\draw[arr] (x86) -- (asm)
  node[lbl, right, midway] {AT\&T asm};

% Rama de verificacion
\node[box, right=2.5cm of opt, fill=green!8, align=center]
  (wpgen)  {VCGenerator\\WPCalculus};
\node[box, below=0.6cm of wpgen, fill=green!8, align=center]
  (smt)    {SMTExporter\\(.smt2)};
\node[box, below=0.6cm of smt, fill=green!8]
  (z3)     {Z3 Solver};

\draw[arr, dashed, green!60!black]
  (sem) -| node[lbl, above, midway] {--verify} (wpgen);
\draw[arr, green!60!black] (wpgen) -- (smt);
\draw[arr, green!60!black] (smt) -- (z3);

\end{tikzpicture}
\caption{Pipeline completo del compilador JavaSubset. La rama verde
  (verificacion formal) se activa con la opcion \texttt{--verify}.}
\label{fig:pipeline}
\end{figure}

\section{Representaciones intermedias}

A lo largo del pipeline el programa adopta cuatro representaciones distintas:

\begin{enumerate}
  \item \textbf{Codigo fuente}: texto Java con posibles anotaciones formales.
  \item \textbf{AST} (\emph{Abstract Syntax Tree}): arbol de nodos tipados,
    jerarquia de clases C++ visitables mediante el patron \emph{Visitor}.
  \item \textbf{Cuadruplas IR}: secuencia plana de instrucciones de la forma
    $(op, \mathit{arg}_1, \mathit{arg}_2, \mathit{res})$, independiente de la
    arquitectura.
  \item \textbf{Ensamblador x86-64}: instrucciones AT\&T con registros reales
    o con posiciones en la pila.
\end{enumerate}

\section{Modos de operacion}

El compilador admite tres modos de operacion seleccionables mediante flags:

\begin{table}[H]
  \centering
  \caption{Flags del compilador y su efecto}
  \label{tab:flags}
  \begin{tabular}{llp{7cm}}
    \toprule
    \textbf{Flag}        & \textbf{Salida}   & \textbf{Descripcion} \\
    \midrule
    (ninguno)            & \texttt{prog.s}   & Compilacion estandar a ensamblador \\
    \texttt{--ir}        & \texttt{stderr}   & Imprime cuadruplas IR y forma lineal \\
    \texttt{--reg-alloc} & \texttt{stderr}   & Imprime mapa de asignacion de registros \\
    \texttt{--verify}    & \texttt{prog.smt2}& Verificacion formal con Z3 \\
    \bottomrule
  \end{tabular}
\end{table}

\section{Patron Visitor}

Todas las fases que procesan el AST utilizan el patron de diseno
\textbf{Visitor} (GoF). La interfaz abstracta \texttt{Visitor} declara un
metodo \texttt{visit()} para cada tipo de nodo del AST; cada fase
implementa esta interfaz sin modificar las clases del arbol.

\begin{figure}[H]
\centering
\begin{tikzpicture}[
  iface/.style={rectangle, draw, fill=yellow!15,
                minimum width=4cm, minimum height=0.7cm,
                align=center, font=\small},
  impl/.style={rectangle, draw, fill=blue!8,
               minimum width=3.5cm, minimum height=0.7cm,
               align=center, font=\small},
  arr/.style={-Stealth, thick}
]
\node[iface] (v) {\texttt{Visitor} (abstract)};
\node[impl, below left=1cm and 2.5cm of v] (sv) {\texttt{SemanticVisitor}};
\node[impl, below left=1cm and 0cm of v]   (ib) {\texttt{IRBuilder}};
\node[impl, below right=1cm and 0cm of v]  (av) {\texttt{ASTVisualizer}};
\node[impl, below right=1cm and 2.5cm of v] (dv) {\texttt{DOTVisualizer}};

\draw[arr] (sv) -- (v);
\draw[arr] (ib) -- (v);
\draw[arr] (av) -- (v);
\draw[arr] (dv) -- (v);
\end{tikzpicture}
\caption{Jerarquia del patron Visitor en el compilador.}
\label{fig:visitor}
\end{figure}

\section{Estructura de directorios}

\begin{lstlisting}[style=pseudo, caption={Organizacion del repositorio}]
.
|-- include/         Cabeceras: ASTNodes.h, IR.h, Predicate.h, ...
|-- src/             Implementaciones: lexer.l, parser.y, *Visitor.h, ...
|   |-- lexer.l           Analizador lexico (Flex)
|   |-- parser.y          Gramatica + acciones semanticas (Bison)
|   |-- SemanticVisitor.h Comprobacion de tipos
|   |-- IRBuilder.h       Traduccion AST -> IR
|   |-- IROptimizer.h     Optimizaciones sobre IR
|   |-- RegisterAllocator.h  Chaitin-Briggs coloring
|   |-- X86Generator.h    Generacion de codigo x86-64
|   |-- WPCalculus.h      Calculo de precondicion mas debil
|   |-- VCGenerator.h     Generacion de condiciones de verificacion
|   `-- SMTExporter.h     Exportacion a SMT-Lib2 y ejecucion de Z3
|-- tests/
|   |-- backend/          Programas de prueba end-to-end
|   |-- semantic/         Pruebas de deteccion de errores
|   |-- verification/     Demos de verificacion formal
|   `-- run_tests.sh      Script de pruebas
`-- Makefile
\end{lstlisting}

% ─────────────────────────────────────────────────────────────────────────────
%  Capítulo 3 – Frontend: Análisis léxico-sintáctico y AST
% ─────────────────────────────────────────────────────────────────────────────
\chapter{Frontend: An\'alisis l\'exico-sint\'actico}
\label{chap:frontend}

\section{An\'alisis l\'exico}

El analizador l\'exico se implementa en \texttt{src/lexer.l} mediante Flex.
Convierte el flujo de caracteres de entrada en una secuencia de \emph{tokens}
que el parser consumir\'a.

\subsection{Categor\'ias de tokens}

\begin{table}[H]
  \centering
  \caption{Categor\'ias de tokens reconocidas por el lexer}
  \label{tab:tokens}
  \begin{tabular}{llp{6cm}}
    \toprule
    \textbf{Categor\'ia} & \textbf{Patr\'on} & \textbf{Ejemplos} \\
    \midrule
    Palabras reservadas & exactas       & \texttt{int}, \texttt{while}, \texttt{return} \\
    Tipos               & exactas       & \texttt{int}, \texttt{double}, \texttt{boolean} \\
    Anotaciones formales& exactas       & \texttt{@pre}, \texttt{@post}, \texttt{@invariant} \\
    Identificadores     & \texttt{[a-zA-Z][a-zA-Z0-9]*} & \texttt{x}, \texttt{sum}, \texttt{absValue} \\
    Literales enteros   & \texttt{[0-9]+}               & \texttt{0}, \texttt{42} \\
    Literales reales    & \texttt{[0-9]+\textbackslash.[0-9]+} & \texttt{3.14} \\
    Operadores          & s\'imbolos    & \texttt{+}, \texttt{-}, \texttt{==}, \texttt{==>} \\
    Puntuaci\'on        & s\'imbolos    & \texttt{\{}, \texttt{;}, \texttt{(} \\
    \bottomrule
  \end{tabular}
\end{table}

\subsection{Tokens especiales para verificaci\'on}

Para la fase 5 se a\~nadieron tokens dedicados a las anotaciones de
especificaci\'on formal, colocados \emph{antes} de la regla general de
identificadores para que tengan precedencia:

\begin{lstlisting}[style=cpp, caption={Reglas l\'exicas para anotaciones formales}]
"result"      { RET_ATTR(result_token); }
"@pre"        { RET_ATTR(pre_token); }
"@post"       { RET_ATTR(post_token); }
"@invariant"  { RET_ATTR(invariant_token); }
"@variant"    { RET_ATTR(variant_token); }
"==>"         { RET_ATTR(implies_token); }
\end{lstlisting}

\section{An\'alisis sint\'actico}

El parser se implementa en \texttt{src/parser.y} con GNU Bison (gram\'atica
LALR(1)). Construye el AST de forma dirigida por la sintaxis mediante
acciones sem\'anticas en C++.

\subsection{Estructura gramatical principal}

\begin{lstlisting}[style=pseudo, caption={Reglas principales de la gram\'atica Bison}]
// Un programa es una clase con miembros
Program : public class id '(' ')' '{' Members '}'
Members : Members Member | Members OptPre | Members OptPost | /* empty */
Member  : Main | Method | DVar

// Metodo estatico sin parametros
Method  : public static Tipo id '(' ')' Bloque
        | public static void id '(' ')' Bloque

// Instrucciones
Instr   : DVar | Asig | Print | Return
        | if '(' Expr ')' Instr
        | if '(' Expr ')' Instr else Instr
        | while '(' Expr ')' OptInvariant OptVariant Instr
        | Bloque

// Anotaciones de especificacion formal
OptPre      : '@pre'       PredExpr | /* empty */
OptPost     : '@post'      PredExpr | /* empty */
OptInvariant: '@invariant' PredExpr | /* empty */
OptVariant  : '@variant'   PredExpr | /* empty */
PredExpr    : Expr '==>' Expr | Expr
\end{lstlisting}

\subsection{Gesti\'on de anotaciones formales}

Las anotaciones \texttt{@pre}/\texttt{@post} se acumulan en variables globales
\texttt{g\_pre} y \texttt{g\_post} antes de que se procese el m\'etodo. Cuando
se reduce una regla \texttt{Method}, el nodo \texttt{MethodNode} recibe las
precondici\'on y postcondici\'on almacenadas:

\begin{lstlisting}[style=cpp, caption={Vinculaci\'on de anotaciones al MethodNode}]
Method : public_token static_token Tipo id pari pard Bloque {
    MethodNode* m = new MethodNode($4->lexema, $3->tipo, $7, ...);
    m->precondition.reset(g_pre);   // @pre almacenada globalmente
    m->postcondition.reset(g_post); // @post almacenada globalmente
    g_pre = g_post = nullptr;
    $$ = m;
}
\end{lstlisting}

\section{\'Arbol Sint\'actico Abstracto (AST)}

El AST se representa mediante una jerarqu\'ia de clases C++ en
\texttt{include/ASTNodes.h}.

\subsection{Jerarqu\'ia de nodos}

\begin{figure}[H]
\centering
\begin{tikzpicture}[
  cls/.style={rectangle, draw, rounded corners=2pt,
              minimum width=3cm, minimum height=0.6cm,
              align=center, font=\small\ttfamily, fill=blue!8},
  arr/.style={-Stealth, thick},
  node distance=0.5cm and 1.2cm
]
\node[cls] (node)   {Node};
\node[cls, below left=0.8cm and 1.5cm of node]  (stmt)  {StmtNode};
\node[cls, below right=0.8cm and 1.5cm of node] (expr)  {ExprNode};

\node[cls, below=0.5cm of stmt] (assign) {AssignNode};
\node[cls, below=0.5cm of assign] (block)  {BlockNode};
\node[cls, below=0.5cm of block]  (ifn)    {IfNode};
\node[cls, below=0.5cm of ifn]    (whi)    {WhileNode};
\node[cls, below=0.5cm of whi]    (met)    {MethodNode};
\node[cls, below=0.5cm of met]    (ret)    {ReturnNode};

\node[cls, below=0.5cm of expr]   (bin)    {BinaryExprNode};
\node[cls, below=0.5cm of bin]    (una)    {UnaryExprNode};
\node[cls, below=0.5cm of una]    (ide)    {IdentifierNode};
\node[cls, below=0.5cm of ide]    (ilit)   {IntLiteralNode};
\node[cls, below=0.5cm of ilit]   (flo)    {FloatLiteralNode};
\node[cls, below=0.5cm of flo]    (cal)    {CallNode};
\node[cls, below=0.5cm of cal]    (res)    {ResultNode};

\draw[arr] (stmt) -- (node);
\draw[arr] (expr) -- (node);
\draw[arr] (assign) -- (stmt);
\draw[arr] (block) -- (stmt);
\draw[arr] (ifn) -- (stmt);
\draw[arr] (whi) -- (stmt);
\draw[arr] (met) -- (stmt);
\draw[arr] (ret) -- (stmt);
\draw[arr] (bin) -- (expr);
\draw[arr] (una) -- (expr);
\draw[arr] (ide) -- (expr);
\draw[arr] (ilit) -- (expr);
\draw[arr] (flo) -- (expr);
\draw[arr] (cal) -- (expr);
\draw[arr] (res) -- (expr);
\end{tikzpicture}
\caption{Jerarqu\'ia de nodos del AST (flecha = hereda de).}
\label{fig:ast}
\end{figure}

\subsection{Nodos especiales para verificaci\'on}

\begin{description}
  \item[\texttt{ResultNode}] Representa la palabra reservada \texttt{result}
    en las postcondiciones. Hace referencia al valor de retorno del m\'etodo sin
    necesidad de nombrarlo expl\'icitamente.
  \item[\texttt{MethodNode.precondition / .postcondition}] Campos
    \texttt{unique\_ptr<ExprNode>} que almacenan las expresiones de las
    anotaciones \texttt{@pre} y \texttt{@post}.
  \item[\texttt{WhileNode.invariant / .variant}] Campos an\'alogos para las
    anotaciones \texttt{@invariant} y \texttt{@variant} del bucle.
\end{description}

\begin{lstlisting}[style=cpp, caption={Nodo WhileNode con soporte de invariante}]
class WhileNode : public StmtNode {
public:
    std::unique_ptr<ExprNode> condition;
    std::unique_ptr<StmtNode> body;
    std::unique_ptr<ExprNode> invariant;  // @invariant annotation
    std::unique_ptr<ExprNode> variant;    // @variant annotation
    ...
};
\end{lstlisting}

% ─────────────────────────────────────────────────────────────────────────────
%  Capítulo 4 – Análisis Semántico
% ─────────────────────────────────────────────────────────────────────────────
\chapter{Análisis semántico}
\label{chap:semantica}

\section{Tabla de símbolos}

La tabla de símbolos (\texttt{TablaSimbolos}) es una estructura encadenada de
ámbitos. Cada \texttt{BlockNode} posee su propia \texttt{TablaSimbolos} apuntando
a la tabla del bloque padre, lo que implementa el alcance léxico estándar de Java.

\begin{lstlisting}[style=cpp, caption={Estructura de la tabla de símbolos}]
struct Simbolo {
    string nombre;
    unsigned tipo;          // ENTERO=0, REAL=1, LOGICO=2, SCVAR=3
    unsigned dir;           // desplazamiento en la pila (para IR/x86)
    unsigned tam;
    bool isConstant;        // declarado con 'final'
    bool isAssigned;        // asignado al menos una vez
    bool isFunction;        // true if this symbol is a method
    int  returnType;        // return type (if isFunction)
};

class TablaSimbolos {
    TablaSimbolos* padre;           // parent scope (lexical scoping)
    vector<Simbolo> simbolos;
    bool set(Simbolo s);            // insert; returns false if duplicate
    Simbolo* get(string nombre);    // lookup in this scope and parent
};
\end{lstlisting}

\section{Comprobaciones semánticas}

El \texttt{SemanticVisitor} recorre el AST tipado e implementa:

\begin{enumerate}
  \item \textbf{Declaración antes de uso}: verifica que toda variable
    referenciada esté declarada en algún ámbito visible.
  \item \textbf{Compatibilidad de tipos}: comprueba que los operandos de
    cada operador sean compatibles (p.ej. no se puede sumar un \texttt{int}
    y un \texttt{boolean}).
  \item \textbf{Constantes}: detecta reasignaciones a variables \texttt{final}.
  \item \textbf{Tipo de retorno}: verifica que las sentencias \texttt{return}
    devuelvan un valor compatible con el tipo de retorno del método.
  \item \textbf{Código inalcanzable}: detecta sentencias tras un
    \texttt{return} incondicional.
  \item \textbf{Llamadas a funciones}: comprueba que el nombre de la función
    esté declarado y tiene el tipo de retorno correcto.
\end{enumerate}

\subsection{Reglas de coerción de tipos}

\begin{table}[H]
  \centering
  \caption{Tabla de coerción de tipos para operaciones binarias}
  \label{tab:coercion}
  \begin{tabular}{cccl}
    \toprule
    \textbf{Operando izq.} & \textbf{Op.} & \textbf{Operando der.} & \textbf{Tipo resultado} \\
    \midrule
    \texttt{int}    & $+,-,*,/,\%$  & \texttt{int}    & \texttt{int} \\
    \texttt{double} & $+,-,*,/$     & \texttt{double} & \texttt{double} \\
    \texttt{int}    & $+,-,*,/$     & \texttt{double} & \texttt{double} (promoución) \\
    cualquiera      & $==,!=,<,>$   & mismo           & \texttt{boolean} \\
    \texttt{boolean}& $\&\&, \|$    & \texttt{boolean}& \texttt{boolean} \\
    \bottomrule
  \end{tabular}
\end{table}

\section{Ejemplo de error semántico detectado}

\begin{lstlisting}[style=java, caption={Reasignación de constante — error semántico}]
public class EjError {
    public static void main(String[] args) {
        final int MAX;
        MAX = 100;
        MAX = 200;  // Error: reasignacion de constante 'MAX'
    }
}
\end{lstlisting}

El compilador emite:
\begin{verbatim}
Error semantico (5,9): reasignacion a constante 'MAX'
\end{verbatim}

\section{Visita a nodos de especificación formal}

El \texttt{SemanticVisitor} también visita los nodos de anotación para
comprobar su coherencia:

\begin{lstlisting}[style=cpp, caption={Visita a ResultNode en el análisis semántico}]
void visit(ResultNode* node) override {
    if (currentReturnType == -1) {
        error("'result' solo puede usarse en postcondiciones "
              "de metodos no void");
        return;
    }
    node->type = currentReturnType;
}
\end{lstlisting}

% ─────────────────────────────────────────────────────────────────────────────
%  Capítulo 5 – Representación Intermedia
% ─────────────────────────────────────────────────────────────────────────────
\chapter{Representación intermedia}
\label{chap:ir}

\section{Motivación}

La representación intermedia (\IR{}) desacopla el front-end
(específico del lenguaje fuente) del back-end (específico de la
arquitectura destino). Esto permite, en principio, añadir nuevos
lenguajes fuente o nuevas arquitecturas destino sin reescribir
el compilador completo.

\section{Formato de cuádruplas}

La \IR{} consiste en una secuencia plana de \textbf{cuádruplas} de la forma:
\[
  (op,\; \mathit{arg}_1,\; \mathit{arg}_2,\; \mathit{result})
\]

\begin{table}[H]
  \centering
  \caption{Opcodes de la representación intermedia}
  \label{tab:opcodes}
  \begin{tabular}{llp{7cm}}
    \toprule
    \textbf{Opcode}   & \textbf{Ejemplo}          & \textbf{Significado} \\
    \midrule
    \texttt{IR\_ADD}  & $(+, a, b, t_0)$         & $t_0 \leftarrow a + b$ \\
    \texttt{IR\_SUB}  & $(-, a, b, t_0)$         & $t_0 \leftarrow a - b$ \\
    \texttt{IR\_MUL}  & $(*, a, b, t_0)$         & $t_0 \leftarrow a * b$ \\
    \texttt{IR\_DIV}  & $(/, a, b, t_0)$         & $t_0 \leftarrow a / b$ \\
    \texttt{IR\_MOD}  & $(\%, a, b, t_0)$        & $t_0 \leftarrow a \bmod b$ \\
    \texttt{IR\_COPY} & $(=, a, -, x)$           & $x \leftarrow a$ \\
    \texttt{IR\_JZ}   & $(jz, cond, L, -)$       & \texttt{if} cond $= 0$ \texttt{goto} $L$ \\
    \texttt{IR\_JMP}  & $(jmp, L, -, -)$         & \texttt{goto} $L$ \\
    \texttt{IR\_LABEL}& $(label, L, -, -)$       & marca de salto $L$ \\
    \texttt{IR\_CALL} & $(call, f, -, t_0)$      & $t_0 \leftarrow f()$ \\
    \texttt{IR\_RETURN}& $(ret, a, -, -)$        & devuelve $a$ \\
    \texttt{IR\_PRINT\_INT}& $(print\_int, a, -, -)$ & imprime entero \\
    \texttt{IR\_READ\_INT} & $(read\_int, -, -, t_0)$ & lee entero de stdin \\
    \texttt{IR\_HALT} & $(halt, -, -, -)$        & fin del programa \\
    \bottomrule
  \end{tabular}
\end{table}

\section{Argumentos IR}

Cada argumento de una cuádrupla es un \texttt{IRArg} que puede ser:

\begin{itemize}
  \item \textbf{Constante entera}: \texttt{IRArg::constInt(42)}
  \item \textbf{Constante real}: \texttt{IRArg::constReal(3.14)}
  \item \textbf{Variable}: \texttt{IRArg::var("x", addr, tipo)}
  \item \textbf{Temporal}: \texttt{prog.newTemp(tipo)} — generado automáticamente
  \item \textbf{Etiqueta}: \texttt{IRArg::label("L0")}
  \item \textbf{Función}: \texttt{IRArg::func("absValue")}
  \item \textbf{Ninguno}: \texttt{IRArg::none()}
\end{itemize}

\section{Traducción de construcciones}

\subsection{Expresión aritmética}

\begin{ejemplo}
La expresión \texttt{(a + b) * c} se traduce como:
\[
  (IR\_ADD, a, b, t_0),\quad (IR\_MUL, t_0, c, t_1)
\]
\end{ejemplo}

\subsection{Sentencia if-else}

\begin{lstlisting}[style=java, caption={if-else en código fuente}]
if (x > 0) y = x; else y = -x;
\end{lstlisting}

\begin{verbatim}
(IR_GT,    x, 0,  t0)      // t0 = x > 0
(IR_JZ,    t0, L_else, -)  // if !t0 goto L_else
(IR_COPY,  x, -, y)        // then: y = x
(IR_JMP,   L_end, -, -)
(IR_LABEL, L_else, -, -)
(IR_NEG,   x, -, t1)       // t1 = -x
(IR_COPY,  t1, -, y)       // else: y = -x
(IR_LABEL, L_end, -, -)
\end{verbatim}

\subsection{Bucle while}

\begin{verbatim}
(IR_LABEL, L_start, -, -)
(IR_LTE,   i, 5, t0)
(IR_JZ,    t0, L_end, -)
... cuerpo ...
(IR_JMP,   L_start, -, -)
(IR_LABEL, L_end, -, -)
\end{verbatim}

\subsection{Llamada a función}

\begin{verbatim}
(IR_FUNC_BEGIN, absValue, -, -)
... cuerpo de absValue ...
(IR_RETURN, t3, -, -)
(IR_FUNC_END, absValue, -, -)

... en el punto de llamada ...
(IR_CALL, absValue, -, t5)  // t5 = absValue()
\end{verbatim}

\section{Visualización de la IR}

Con el flag \texttt{--ir} el compilador imprime las cuádruplas en formato
legible por el programador:

\begin{verbatim}
$ ./compiler --ir tests/backend/t04_while.java 2>&1 | head -20
[FUNC_BEGIN] sumToFive
[LABEL] L0
[LTE] i 5 -> t0
[JZ] t0 -> L1
[ADD] s i -> t1
[COPY] t1 -> s
[ADD] i 1 -> t2
[COPY] t2 -> i
[JMP] L0
[LABEL] L1
[RETURN] s
[FUNC_END] sumToFive
\end{verbatim}

% ─────────────────────────────────────────────────────────────────────────────
%  Capítulo 6 – Backend x86-64
% ─────────────────────────────────────────────────────────────────────────────
\chapter{Backend x86-64}
\label{chap:backend}

\section{Convención de llamada System V AMD64 ABI}

El compilador respeta la convención de llamada estándar de Linux/macOS
para x86-64 (\emph{System V AMD64 ABI}):

\begin{table}[H]
  \centering
  \caption{Registros de la System V AMD64 ABI}
  \label{tab:abi}
  \begin{tabular}{lll}
    \toprule
    \textbf{Propósito}     & \textbf{Enteros}   & \textbf{Flotantes} \\
    \midrule
    Argumentos             & \reg{rdi}, \reg{rsi}, \reg{rdx}, \reg{rcx}, \reg{r8}, \reg{r9}
                           & \reg{xmm0}–\reg{xmm7} \\
    Valor de retorno       & \reg{rax}           & \reg{xmm0} \\
    Callee-saved           & \reg{rbx}, \reg{r12}–\reg{r15}, \reg{rbp} & \reg{xmm8}–\reg{xmm15} \\
    Caller-saved           & \reg{rcx}, \reg{rdx}, \reg{rsi}, \reg{rdi}, \reg{r8}–\reg{r11}
                           & \reg{xmm0}–\reg{xmm7} \\
    \bottomrule
  \end{tabular}
\end{table}

\section{Marco de pila}

Cada función establece su propio marco de pila siguiendo el
prólogo/epílogo estándar:

\begin{lstlisting}[style=asm, caption={Prólogo y epílogo de función}]
    pushq   %rbp
    movq    %rsp, %rbp
    subq    $N, %rsp        # N = frameSize (alineado a 16 bytes)
    ...
    movq    %rbp, %rsp
    popq    %rbp
    ret
\end{lstlisting}

Las variables locales y temporales se ubican en la pila según la fórmula:
\[
  \texttt{slot}(v) = -((\text{slot}_v + 1) \times 8)\texttt{(\%rbp)}
\]
donde $\text{slot}_v$ es el índice (0-based) asignado a la primera vez
que se usa la dirección de la variable en la IR.

\section{Generación de código por opcode}

\subsection{Operaciones aritméticas enteras}

\begin{table}[H]
  \centering
  \caption{Traducción de operaciones aritméticas enteras}
  \label{tab:arith}
  \begin{tabular}{lp{8cm}}
    \toprule
    \textbf{IR}     & \textbf{Secuencia x86-64 AT\&T} \\
    \midrule
    $t \leftarrow a + b$ &
      \texttt{movq mem(a), \%rax; addq mem(b), \%rax; movq \%rax, mem(t)} \\
    $t \leftarrow a * b$ &
      \texttt{movq mem(a), \%rax; imulq mem(b); movq \%rax, mem(t)} \\
    $t \leftarrow a / b$ &
      \texttt{movq mem(a), \%rax; cqto; idivq mem(b); movq \%rax, mem(t)} \\
    $t \leftarrow a \% b$ &
      \texttt{movq mem(a), \%rax; cqto; idivq mem(b); movq \%rdx, mem(t)} \\
    \bottomrule
  \end{tabular}
\end{table}

\subsection{Saltos condicionales}

Los saltos se implementan con el par \texttt{cmpq} + instrucción de salto
condicional. Para comparaciones que producen un booleano en la IR (p.ej.
\texttt{IR\_LT}), el patrón es:

\begin{lstlisting}[style=asm, caption={Comparación IR\_LT y salto condicional}]
    movq    mem(a), %rax
    cmpq    mem(b), %rax
    setl    %al              # 1 si a < b, 0 si no
    movzbq  %al, %rax
    movq    %rax, mem(t)
    ...
    testq   mem(cond), mem(cond)
    je      L_else           # JZ: salta si cond == 0
\end{lstlisting}

\subsection{Llamadas a función}

\begin{lstlisting}[style=asm, caption={Llamada a función y recogida de valor de retorno}]
    call    absValue         # llama a absValue()
    movq    %rax, mem(t)    # t = valor de retorno (%rax)
\end{lstlisting}

\subsection{E/S con \texttt{printf} y \texttt{scanf}}

\begin{lstlisting}[style=asm, caption={Impresión de entero con printf}]
    .section .rodata
fmt_int:    .string "%ld\n"
    .text
    leaq    fmt_int(%rip), %rdi
    movq    mem(val), %rsi
    xorl    %eax, %eax
    call    printf@PLT
\end{lstlisting}

\begin{lstlisting}[style=asm, caption={Lectura de entero con scanf}]
fmt_scanf_int: .string "%ld"
    leaq    fmt_scanf_int(%rip), %rdi
    leaq    mem(t)(%rbp), %rsi       # address of local variable in stack
    xorl    %eax, %eax
    call    scanf@PLT
\end{lstlisting}

\section{Ejemplo completo de programa compilado}

\begin{lstlisting}[style=java, caption={Programa fuente: suma hasta 5}]
public class TestWhile {
    public static int sumToFive() {
        int i; int s;
        i = 1; s = 0;
        while (i <= 5) { s = s + i; i = i + 1; }
        return s;
    }
    public static void main(String[] args) {
        int val;
        val = sumToFive();
        System.out.println(val);
    }
}
\end{lstlisting}

El ensamblador generado para \texttt{sumToFive} tiene la forma:

\begin{lstlisting}[style=asm, caption={Ensamblador x86-64 generado para sumToFive}]
    .globl  sumToFive
sumToFive:
    pushq   %rbp
    movq    %rsp, %rbp
    subq    $48, %rsp
L0:
    movq    -8(%rbp), %rax      # i
    cmpq    $5, %rax
    setle   %al
    movzbq  %al, %rax
    movq    %rax, -24(%rbp)     # t0 = i <= 5
    testq   -24(%rbp), %rax
    je      L1
    movq    -8(%rbp), %rax      # s = s + i
    addq    -16(%rbp), %rax
    movq    %rax, -32(%rbp)
    movq    -32(%rbp), %rax
    movq    %rax, -16(%rbp)
    movq    -8(%rbp), %rax      # i = i + 1
    addq    $1, %rax
    ...
    jmp     L0
L1:
    movq    -16(%rbp), %rax     # return s
    movq    %rbp, %rsp
    popq    %rbp
    ret
\end{lstlisting}

% ─────────────────────────────────────────────────────────────────────────────
%  Capítulo 7 – Optimizaciones sobre la IR
% ─────────────────────────────────────────────────────────────────────────────
\chapter{Optimizaciones sobre la representación intermedia}
\label{chap:optimizacion}

\section{Introducción}

Las optimizaciones se aplican sobre la \IR{} antes de la generación de
código. Al operar sobre las cuádruplas —y no sobre el ensamblador— las
optimizaciones son independientes de la arquitectura destino.

El módulo \texttt{IROptimizer} implementa las siguientes pasadas:

\begin{enumerate}
  \item Eliminación de código muerto (\emph{dead code elimination}).
  \item Plegado de constantes (\emph{constant folding}).
  \item Propagación de copias (\emph{copy propagation}).
\end{enumerate}

\section{Eliminación de código muerto}

Se eliminan cuádruplas que calculan un valor que nunca se usa
(temporales sin ningún uso posterior) y cuádruplas que son
literalmente inalcanzables (tras un \texttt{IR\_JMP} o \texttt{IR\_RETURN}
incondicional antes de una etiqueta).

\begin{definicion}[Código muerto]
  Una cuádrupla $q = (op, a_1, a_2, t)$ es \emph{código muerto} si
  $t$ no es una variable de usuario, no aparece como argumento en
  ninguna cuádrupla posterior, y $op \notin \{IR\_CALL, IR\_READ\_INT,
  IR\_READ\_REAL, IR\_PRINT\_*, IR\_JMP, IR\_JZ, IR\_LABEL, IR\_RETURN,
  IR\_HALT\}$ (instrucciones con efecto de borde).
\end{definicion}

La eliminación se realiza en dos pasadas:

\begin{enumerate}
  \item Calcular el conjunto $\textit{used}$ de todas las direcciones que
    aparecen como argumento en alguna cuádrupla.
  \item Eliminar las cuádruplas cuyo resultado no está en $\textit{used}$
    y que no tienen efecto de borde.
\end{enumerate}

\section{Plegado de constantes}

Si ambos operandos de una operación aritmética o de comparación son constantes
literales, se puede calcular el resultado en tiempo de compilación:

\begin{ejemplo}
  La cuádrupla $(IR\_ADD, 3, 5, t_0)$ se reemplaza directamente por
  $t_0 \leftarrow 8$, es decir, una cuádrupla $(IR\_COPY, 8, -, t_0)$.
\end{ejemplo}

\begin{table}[H]
  \centering
  \caption{Operaciones que admiten plegado de constantes}
  \label{tab:folding}
  \begin{tabular}{ll}
    \toprule
    \textbf{Operación} & \textbf{Resultado} \\
    \midrule
    $\text{int} + \text{int}$   & \texttt{constInt(a+b)} \\
    $\text{int} - \text{int}$   & \texttt{constInt(a-b)} \\
    $\text{int} * \text{int}$   & \texttt{constInt(a*b)} \\
    $\text{int} / \text{int}$, $b \neq 0$ & \texttt{constInt(a/b)} \\
    $\text{double} \oplus \text{double}$ & \texttt{constReal(a $\oplus$ b)} \\
    $a == a$, $a \leq a$        & \texttt{constInt(1)} \\
    $a < a$, $a \neq a$         & \texttt{constInt(0)} \\
    \bottomrule
  \end{tabular}
\end{table}

\section{Propagación de copias}

Si una cuádrupla tiene la forma $(IR\_COPY, a, -, x)$ y $a$ es un
argumento constante o una variable que no se modifica en el tramo
hasta el uso siguiente de $x$, se puede sustituir $x$ por $a$
directamente, eliminando la copia.

\begin{ejemplo}
  La secuencia:
  \begin{verbatim}
  (IR_COPY, 42, -, t0)
  (IR_ADD,  t0,  5, t1)
  \end{verbatim}
  se optimiza a:
  \begin{verbatim}
  (IR_ADD, 42, 5, t1)   ->  (IR_COPY, 47, -, t1)  [plegado posterior]
  \end{verbatim}
\end{ejemplo}

\section{Impacto de las optimizaciones}

\begin{table}[H]
  \centering
  \caption{Reducción de cuádruplas en los programas de prueba}
  \label{tab:opt-impact}
  \begin{tabular}{lrrr}
    \toprule
    \textbf{Programa} & \textbf{Sin opt.} & \textbf{Con opt.} & \textbf{Reducción} \\
    \midrule
    \texttt{t01\_print.java}      & 8  & 5  & $-37.5\%$ \\
    \texttt{t02\_arithmetic.java} & 22 & 14 & $-36.4\%$ \\
    \texttt{t03\_if\_else.java}   & 18 & 12 & $-33.3\%$ \\
    \texttt{t04\_while.java}      & 26 & 20 & $-23.1\%$ \\
    \texttt{t05\_functions.java}  & 15 & 11 & $-26.7\%$ \\
    \bottomrule
  \end{tabular}
\end{table}

% ─────────────────────────────────────────────────────────────────────────────
%  Capítulo 8 – Asignación de Registros (Chaitin-Briggs)
% ─────────────────────────────────────────────────────────────────────────────
\chapter{Asignación de registros: Chaitin-Briggs}
\label{chap:registro}

\section{Motivación}

Los registros de la CPU son escasos y de acceso mucho más rápido que la memoria.
La \emph{asignación de registros} consiste en asignar las variables y temporales
de la \IR{} a registros físicos (o a posiciones en la pila si no hay
registros suficientes: \emph{spilling}).

El algoritmo de \textbf{Chaitin-Briggs}~\cite{chaitin1982} modela el
problema como un \emph{problema de coloración de grafos}: dos variables
que están \emph{vivas} en el mismo punto del programa no pueden compartir
el mismo registro, y por tanto se conectan en el \emph{grafo de interferencias}.
Colorear el grafo con $K$ colores equivale a asignar $K$ registros.

\section{Registros disponibles}

Se reservan los siguientes registros \emph{caller-saved} para temporales,
evitando tener que salvarlos en llamadas a función:

\begin{table}[H]
  \centering
  \caption{Registros asignables por el compilador}
  \label{tab:regs}
  \begin{tabular}{ll}
    \toprule
    \textbf{Clase} & \textbf{Registros} ($K = 6$) \\
    \midrule
    Enteros  & \reg{rcx}, \reg{rdx}, \reg{r8}, \reg{r9}, \reg{r10}, \reg{r11} \\
    Flotantes & \reg{xmm2}, \reg{xmm3}, \reg{xmm4}, \reg{xmm5}, \reg{xmm6}, \reg{xmm7} \\
    \bottomrule
  \end{tabular}
\end{table}

\section{Algoritmo}

El algoritmo consta de tres fases:

\subsection{Fase 1: Análisis de vivacidad}

Se calculan los conjuntos \emph{use} y \emph{def} por bloque básico y
se resuelven las ecuaciones de punto fijo:
\[
  \text{live\_in}[B]  = \text{use}[B]  \cup (\text{live\_out}[B] \setminus \text{def}[B])
\]
\[
  \text{live\_out}[B] = \bigcup_{S \in \text{succ}(B)} \text{live\_in}[S]
\]

\begin{lstlisting}[style=pseudo, caption={Análisis de vivacidad (punto fijo)}]
Entrada: bloques B_1,...,B_n con conjuntos use[B], def[B]
Salida:  live_in[B], live_out[B] para cada bloque

repetir hasta que no haya cambios:
    para cada bloque B (en orden inverso):
        live_out[B] = union{ live_in[S] | S sucesor de B }
        live_in[B]  = use[B] union (live_out[B] menos def[B])
\end{lstlisting}

\subsection{Fase 2: Grafo de interferencias}

Se recorre cada bloque básico de forma inversa manteniendo el conjunto
$\text{live}$ de variables vivas en el punto actual. Al definir una
variable $d$ se añade una arista $(d, v)$ para cada $v \in \text{live}$.

\begin{lstlisting}[style=cpp, caption={Construcción del grafo de interferencias}]
for (auto it = block.rbegin(); it != block.rend(); ++it) {
    const auto& q = *it;
    if (q.result.kind == IRArg::TEMP || q.result.kind == IRArg::VAR) {
        int d = q.result.addr;
        for (int live_var : live)
            graph.addEdge(d, live_var);  // d interfiere con live_var
        live.erase(d);                   // d deja de estar vivo antes de def
    }
    // Add uses to live set
    if (q.arg1 es variable/temporal) live.insert(q.arg1.addr);
    if (q.arg2 es variable/temporal) live.insert(q.arg2.addr);
}
\end{lstlisting}

\subsection{Fase 3: Coloración del grafo}

Se sigue el algoritmo \emph{simplify}:

\begin{enumerate}
  \item Mientras exista un nodo con grado $< K$: apilarlo y eliminarlo del grafo.
  \item Si todos los nodos tienen grado $\geq K$: elegir el de mayor grado para
    \emph{spill} y eliminarlo (se marca como derramado).
  \item Repetir hasta vaciar el grafo.
  \item Desapilar asignando colores (registros) no usados por vecinos ya coloreados.
\end{enumerate}

\begin{lstlisting}[style=cpp, caption={Coloración simplificada de Chaitin-Briggs}]
// Fase simplify: apila nodos con grado < K
stack<int> nodeStack;
while (!worklist.empty()) {
    auto it = find_if(worklist, [&](int n){ return degree[n] < K; });
    if (it != worklist.end()) {
        nodeStack.push(*it);
        removeNode(*it);
    } else {
        // spill: eliminar nodo de mayor grado
        int spill = *max_element(worklist, [&](int a, int b){
            return degree[a] < degree[b]; });
        spilledNodes.insert(spill);
        removeNode(spill);
    }
}
// Fase assign: asignar colores
while (!nodeStack.empty()) {
    int n = nodeStack.top(); nodeStack.pop();
    set<string> usedColors;
    for (int neighbor : adj[n])
        if (color.count(neighbor)) usedColors.insert(color[neighbor]);
    for (const auto& reg : AVAILABLE_REGS)
        if (!usedColors.count(reg)) { color[n] = reg; break; }
}
\end{lstlisting}

\section{Uso del resultado}

El mapa de coloración \texttt{RegMap} (dirección IR → nombre de registro)
se pasa al \texttt{X86Generator}. Las variables con \texttt{color = ""}
(derramadas) se mantienen en la pila y se acceden con la notación
\texttt{-N(\%rbp)}.

\begin{verbatim}
$ ./compiler --reg-alloc tests/backend/t04_while.java 2>&1
=== Asignacion de Registros ===
  addr 1 (i) -> %rcx
  addr 2 (s) -> %rdx
  addr 3 (t0) -> %r8
  addr 4 (t1) -> spill
\end{verbatim}

% ─────────────────────────────────────────────────────────────────────────────
%  Capítulo 9 – Verificación Formal: Lógica de Hoare y Cálculo WP
% ─────────────────────────────────────────────────────────────────────────────
\chapter{Verificación formal: Lógica de Hoare y cálculo WP}
\label{chap:verificacion}

\section{Fundamentos teóricos}

\subsection{Lógica de Hoare}

La \textbf{Lógica de Hoare}~\cite{hoare1969} es un sistema formal para
razonar sobre la corrección de programas imperativos. Introduce las
\emph{tripletas de Hoare}:
\[
  \hoare{P}{S}{Q}
\]
que se leen: ``si la precondición $P$ vale antes de ejecutar $S$,
entonces la postcondición $Q$ vale tras la ejecución de $S$
(asumiendo que $S$ termina)''.

\begin{definicion}[Corrección parcial]
  $\hoare{P}{S}{Q}$ es \emph{parcialmente correcta} si para toda
  ejecución de $S$ que parte de un estado en el que $P$ vale y
  que termina, el estado final satisface $Q$.
\end{definicion}

\begin{definicion}[Corrección total]
  $\hoare{P}{S}{Q}$ es \emph{totalmente correcta} si además
  toda ejecución de $S$ que parte de un estado en el que $P$ vale
  termina necesariamente.
\end{definicion}

\subsection{Cálculo de Precondiciones Más Débiles}

\textbf{Dijkstra}~\cite{dijkstra1976} introduce el \emph{transformador
de predicados} $\WP{S}{Q}$: la \emph{precondición más débil} de $S$
con respecto a $Q$, es decir, el predicado más débil (el que restringe
menos el estado inicial) tal que:
\[
  \hoare{\WP{S}{Q}}{S}{Q}
\]

La condición de verificación principal para un método anotado es:
\[
  \mathit{VC} \;:\; P \Rightarrow \WP{\mathit{body}}{Q}
\]

Si esta fórmula es \textbf{válida} (lo comprueba Z3 como UNSAT de su negación),
el método es parcialmente correcto.

\section{Reglas del cálculo WP}

\subsection{Regla de la asignación (VER-4)}

\begin{regla}[Asignación]
\[
  \WP{x := e}{Q} \;=\; \subst{Q}{x}{e}
\]
donde $\subst{Q}{x}{e}$ denota $Q$ con cada aparición libre de $x$
sustituida por $e$.
\end{regla}

\begin{ejemplo}
  $\WP{x := x+1}{x > 5} = (x+1) > 5$
\end{ejemplo}

\subsection{Regla de la composición (VER-5)}

\begin{regla}[Composición]
\[
  \WP{S_1;\,S_2}{Q} \;=\; \WP{S_1}{\WP{S_2}{Q}}
\]
\end{regla}

Esto implica procesar la secuencia de sentencias en \emph{orden inverso},
comenzando por la postcondición $Q$ y propagando la precondición hacia
atrás.

\paragraph{Retornos anticipados.}
La regla de composición naive es incorrecta en presencia de \texttt{return}
dentro de ramas \texttt{if}. Para la secuencia:
\[
  \texttt{if}(B)\;\texttt{return}\;e_1;\;\;\texttt{return}\;e_2;
\]
la regla correcta es:
\[
  \WP{\cdot}{Q} \;=\; (B \Rightarrow \subst{Q}{\mathtt{result}}{e_1})
                \;\wedge\; (\neg B \Rightarrow \subst{Q}{\mathtt{result}}{e_2})
\]

La implementación usa \textbf{continuaciones} (CPS): cuando se encuentra
un \texttt{return}, se usa $Q$ directamente (no la continuación de la
composición):

\begin{lstlisting}[style=cpp, caption={wpSeq con continuaciones para retornos anticipados}]
static Predicate* wpSeq(const vector<const StmtNode*>& stmts,
                        int idx, const Predicate* Q) {
    if (idx >= stmts.size()) return Q->clone();
    const StmtNode* s = stmts[idx];
    // return: termina el camino, usa Q directamente
    if (auto* r = dynamic_cast<const ReturnNode*>(s))
        return wpReturn(r, Q);
    // if: bifurca la continuacion
    if (auto* i = dynamic_cast<const IfNode*>(s)) {
        Predicate* B = fromExpr(i->condition.get());
        vector<const StmtNode*> cont(stmts.begin()+idx+1, stmts.end());
        Predicate* wp1 = wpBranchAndCont(i->thenPart.get(), cont, Q);
        Predicate* wp2 = wpBranchAndCont(i->elsePart.get(), cont, Q);
        return new AndPred(new ImpliesPred(B->clone(), wp1),
                           new ImpliesPred(new NotPred(B),  wp2));
    }
    // asignacion: WP(x:=e; resto, Q) = WP(resto, Q)[e/x]
    if (auto* a = dynamic_cast<const AssignNode*>(s)) {
        Predicate* restWP = wpSeq(stmts, idx+1, Q);
        return wpAssign(a, restWP);
    }
    return wpSeq(stmts, idx+1, Q);  // otras: sin efecto logico
}
\end{lstlisting}

\subsection{Regla del condicional}

\begin{regla}[Condicional]
\[
  \WP{\texttt{if}(B)\;S_1\;\texttt{else}\;S_2}{Q}
  \;=\;
  (B \Rightarrow \WP{S_1}{Q}) \;\wedge\; (\neg B \Rightarrow \WP{S_2}{Q})
\]
\end{regla}

\subsection{Regla del retorno}

\begin{regla}[Return]
\[
  \WP{\texttt{return}\;e}{Q} \;=\; \subst{Q}{\mathtt{result}}{e}
\]
\end{regla}

\subsection{Regla del bucle con invariante (VER-6)}

Para la verificación de bucles se requiere una \emph{anotación de invariante}
$I$ suministrada por el programador:

\begin{regla}[Bucle con invariante]
  Para $\texttt{while}(B)\;\{S\}$ con invariante $I$ y postcondición $Q$,
  se generan tres \emph{condiciones de verificación}:
  \begin{align*}
    \mathit{VC}_1 \;(\text{iniciación})\; &: P \Rightarrow \WP{S_{\text{antes}}}{I} \\
    \mathit{VC}_2 \;(\text{preservación})\; &: (I \wedge B) \Rightarrow \WP{S_{\text{body}}}{I} \\
    \mathit{VC}_3 \;(\text{uso})\; &: (I \wedge \neg B) \Rightarrow \WP{S_{\text{después}}}{Q}
  \end{align*}
  y el WP del bucle en sí es $I$ (asumiendo la corrección de $I$).
\end{regla}

La precondición más fuerte al entrar en el bucle
($\WP{S_{\text{antes}}}{I}$) se computa correctamente aplicando
\texttt{wpSeq} a las sentencias que preceden al bucle dentro del bloque.

\section{Jerarquía de predicados}

Los predicados lógicos se representan mediante la clase abstracta
\texttt{Predicate} y sus subclases:

\begin{table}[H]
  \centering
  \caption{Clases de predicados implementadas}
  \label{tab:preds}
  \begin{tabular}{llll}
    \toprule
    \textbf{Clase}    & \textbf{toString()}      & \textbf{toSMT()} & \textbf{Uso} \\
    \midrule
    \texttt{TruePred}     & \texttt{true}        & \texttt{true}   & Postcondición trivial \\
    \texttt{FalsePred}    & \texttt{false}       & \texttt{false}  & VC insatisfacible \\
    \texttt{VarPred(x)}   & \texttt{x}           & \texttt{x}      & Variable de programa \\
    \texttt{ResultPred}   & \texttt{result}      & \texttt{result} & Valor de retorno en @post \\
    \texttt{IntLitPred(n)}& $n$                  & $n$             & Constante entera \\
    \texttt{BinaryPred(op,l,r)} & \texttt{(l op r)} & \texttt{(op l r)} & Aritmética/comparación \\
    \texttt{AndPred(l,r)} & \texttt{(l \&\& r)}  & \texttt{(and l r)} & Conjunción \\
    \texttt{OrPred(l,r)}  & \texttt{(l || r)}    & \texttt{(or l r)}  & Disyunción \\
    \texttt{NotPred(p)}   & \texttt{(!p)}        & \texttt{(not p)}   & Negación \\
    \texttt{ImpliesPred(a,c)} & \texttt{(a ==> c)} & \texttt{(=> a c)} & Implicación \\
    \texttt{NegPred(p)}   & \texttt{(-p)}        & \texttt{(- p)}  & Negación aritmética \\
    \bottomrule
  \end{tabular}
\end{table}

La operación clave es la \textbf{sustitución textual}:
$\texttt{subst}(x, e)$ devuelve una copia del predicado donde cada
\texttt{VarPred(x)} se reemplaza por $e$:

\begin{lstlisting}[style=cpp, caption={Sustitución en BinaryPred}]
Predicate* BinaryPred::subst(const string& var, const Predicate* e) const {
    return new BinaryPred(op, left->subst(var, e), right->subst(var, e));
}
Predicate* VarPred::subst(const string& var, const Predicate* e) const {
    if (name == var) return e->clone();  // sustituir
    return new VarPred(name);           // no afectada
}
\end{lstlisting}

\section{Generación de condiciones de verificación (VER-7)}

El \texttt{VCGenerator} genera cuatro tipos de VCs:

\begin{enumerate}
  \item \textbf{Corrección del método}: $P \Rightarrow \WP{\mathit{body}}{Q}$
  \item \textbf{Invariante de bucle} (iniciación, preservación, uso)
  \item \textbf{Seguridad de memoria} (VER-9): para cada división,
    $P \Rightarrow \mathit{divisor} \neq 0$
  \item \textbf{Terminación} (VER-11): para cada bucle con \texttt{@variant} $V$,
    $(I \wedge B) \Rightarrow V \geq 0$
\end{enumerate}

\section{Exportación a SMT-Lib2 (VER-8)}

Cada VC se exporta como un bloque SMT-Lib2. La negación de la VC se envía
a Z3: si Z3 responde \textbf{UNSAT}, la negación es insatisfacible, por lo
tanto la VC es \textbf{válida} (el programa es correcto en ese punto).

\begin{lstlisting}[style=smt, caption={Formato SMT-Lib2 para una VC}]
(push 1)
; [absValue] correctitud
(declare-fun x () Int)     ; variables locales del metodo
(assert (not              ; negacion de la VC
    (=> true              ; precondicion P
        (and
            (=> (< (- 5) 0) (>= (- (- 5)) 0))
            (=> (not (< (- 5) 0)) (>= (- 5) 0))
        )
    )
))
(check-sat)               ; Z3 responde UNSAT = VC valida
(pop 1)
\end{lstlisting}

\section{Interpretación de resultados (VER-10)}

\begin{table}[H]
  \centering
  \caption{Interpretación de la respuesta de Z3}
  \label{tab:z3resp}
  \begin{tabular}{lll}
    \toprule
    \textbf{Respuesta Z3} & \textbf{Significado}         & \textbf{Reporte} \\
    \midrule
    \texttt{unsat}    & VC válida — programa correcto      & \textcolor{green!50!black}{[OK]} \\
    \texttt{sat}      & VC inválida — contraejemplo existe & \textcolor{red!80!black}{[BUG]} \\
    \texttt{unknown}  & Timeout o lógica demasiado compleja & \textcolor{orange}{[?]} \\
    \bottomrule
  \end{tabular}
\end{table}

\section{Ejemplo completo de verificación}

\subsection{Función \texttt{absValue}}

\begin{lstlisting}[style=java, caption={Método verificado: valor absoluto}]
@pre  true
@post result >= 0
public static int absValue() {
    int x;
    x = -5;
    if (x < 0) return -x;
    return x;
}
\end{lstlisting}

Cálculo WP (con continuaciones):

\begin{align*}
  &\WP{\texttt{return x}}{Q_0} = Q_0[\texttt{x}/\texttt{result}] = x \geq 0 \\
  &\WP{\texttt{if}(x<0)\;\texttt{return}\;{-x};\;\texttt{return x}}{Q_0} \\
  &\quad = (x < 0 \Rightarrow Q_0[{-x}/\texttt{result}])
         \wedge (x \geq 0 \Rightarrow Q_0[\texttt{x}/\texttt{result}]) \\
  &\quad = (x < 0 \Rightarrow -x \geq 0) \wedge (x \geq 0 \Rightarrow x \geq 0) \\
  &\WP{x := -5}{\ldots} = [\text{sustituyendo } x = -5] \\
  &\quad = ((-5) < 0 \Rightarrow 5 \geq 0) \wedge ((-5) \geq 0 \Rightarrow (-5) \geq 0) \\
  &\quad = (\top \Rightarrow \top) \wedge (\bot \Rightarrow \bot) = \top
\end{align*}

VC principal: $\true \Rightarrow \true$ — Z3 responde \texttt{unsat} $\Rightarrow$ \textbf{[OK]}.

\subsection{Función \texttt{wrongMethod} (bug intencional)}

\begin{lstlisting}[style=java, caption={Método con bug intencional}]
@post result > 10
public static int wrongMethod() {
    int x;
    x = 5;
    return x;
}
\end{lstlisting}

$\WP{\texttt{return x}}{Q_1} = Q_1[\texttt{x}/\texttt{result}] = x > 10$

$\WP{x := 5}{x > 10} = 5 > 10 = \false$

VC: $\true \Rightarrow \false$ — Z3 responde \texttt{sat} con modelo $x=5 \Rightarrow$ \textbf{[BUG]}.

\subsection{Suma con invariante de bucle}

\begin{lstlisting}[style=java, caption={Suma 1..5 con invariante completo}]
@post result >= 0
public static int sumLoop() {
    int i; int s;
    i = 1; s = 0;
    while (i <= 5)
    @invariant s >= 0 && i >= 1
    @variant   6 - i
    {
        s = s + i;
        i = i + 1;
    }
    return s;
}
\end{lstlisting}

\begin{table}[H]
  \centering
  \caption{VCs generadas para \texttt{sumLoop} y resultados de Z3}
  \label{tab:sumloop-vcs}
  \begin{tabular}{lp{7.5cm}l}
    \toprule
    \textbf{VC} & \textbf{Fórmula} & \textbf{Z3} \\
    \midrule
    Correctitud &
      $\true \Rightarrow (0\geq 0 \wedge 1\geq 1)$ &
      UNSAT \textcolor{green!60!black}{[OK]} \\
    Iniciación &
      $\true \Rightarrow (0\geq 0 \wedge 1\geq 1)$ &
      UNSAT \textcolor{green!60!black}{[OK]} \\
    Preservación &
      $(s\geq 0 \wedge i\geq 1 \wedge i\leq 5) \Rightarrow ((s+i)\geq 0 \wedge (i+1)\geq 1)$ &
      UNSAT \textcolor{green!60!black}{[OK]} \\
    Uso &
      $(s\geq 0 \wedge i\geq 1 \wedge i>5) \Rightarrow s\geq 0$ &
      UNSAT \textcolor{green!60!black}{[OK]} \\
    Terminación &
      $(s\geq 0 \wedge i\geq 1 \wedge i\leq 5) \Rightarrow (6-i)\geq 0$ &
      UNSAT \textcolor{green!60!black}{[OK]} \\
    \bottomrule
  \end{tabular}
\end{table}

\section{Nota sobre invariantes}

Un invariante debe ser \emph{suficientemente fuerte} para que las tres VCs sean
válidas. Un invariante débil como $s \geq 0$ no captura que $i \geq 1$, lo que
hace que Z3 encuentre contraejemplos con $i$ negativo para la VC de
preservación. La fortaleza del invariante es responsabilidad del programador;
el verificador comprueba que el invariante dado es \emph{correcto} (no que sea
\emph{completo}).

% ─────────────────────────────────────────────────────────────────────────────
%  Capítulo 10 – Suite de pruebas
% ─────────────────────────────────────────────────────────────────────────────
\chapter{Suite de pruebas}
\label{chap:pruebas}

\section{Organización}

Las pruebas se organizan en tres grupos principales:

\begin{table}[H]
  \centering
  \caption{Grupos de pruebas}
  \label{tab:grupos}
  \begin{tabular}{llr}
    \toprule
    \textbf{Directorio}            & \textbf{Propósito}                        & \textbf{Tests} \\
    \midrule
    \texttt{tests/backend/}        & Compilación y ejecución completa           & 5 \\
    \texttt{tests/semantic/}       & Errores semánticos esperados               & 3 \\
    \texttt{tests/verification/}   & Verificación formal con Z3                 & 3 \\
    \midrule
    Total                          &                                            & 11 \\
    \bottomrule
  \end{tabular}
\end{table}

El script \texttt{tests/run\_tests.sh} ejecuta todos los grupos y
reporta el resultado de cada prueba como \texttt{[PASS]} o \texttt{[FAIL]}.

\section{Pruebas de backend}

Cada prueba compila el fichero \texttt{.java}, ejecuta el binario
resultante y compara la salida con el valor esperado.

\begin{table}[H]
  \centering
  \caption{Pruebas de backend (compilación y ejecución)}
  \label{tab:backend}
  \begin{tabular}{lp{6cm}l}
    \toprule
    \textbf{Fichero}               & \textbf{Qué se verifica}                  & \textbf{Estado} \\
    \midrule
    \texttt{t01\_print.java}       & \texttt{System.out.println} de enteros     & PASS \\
    \texttt{t02\_arithmetic.java}  & Las cuatro operaciones, módulo y precedencia & PASS \\
    \texttt{t03\_if\_else.java}    & Condicionales con ramas vacías y anidadas  & PASS \\
    \texttt{t04\_while.java}       & Bucles, suma acumulativa, retorno          & PASS \\
    \texttt{t05\_functions.java}   & Llamadas entre métodos, paso de valores    & PASS \\
    \bottomrule
  \end{tabular}
\end{table}

\section{Pruebas semánticas}

Estos ficheros son programas \emph{incorrectos}: el compilador debe
detectar el error semántico y terminar con código de salida distinto
de cero.

\begin{table}[H]
  \centering
  \caption{Pruebas de errores semánticos}
  \label{tab:semantic}
  \begin{tabular}{lp{6.5cm}l}
    \toprule
    \textbf{Fichero}               & \textbf{Error esperado}                   & \textbf{Estado} \\
    \midrule
    \texttt{sem01\_undeclared.java} & Variable usada sin declarar              & PASS \\
    \texttt{sem02\_type\_error.java}& Asignación de \texttt{double} a \texttt{int} sin conversión & PASS \\
    \texttt{sem03\_redecl.java}     & Redeclaración de variable en el mismo ámbito & PASS \\
    \bottomrule
  \end{tabular}
\end{table}

\section{Pruebas de verificación formal}

\begin{table}[H]
  \centering
  \caption{Pruebas de verificación con Z3}
  \label{tab:verif}
  \begin{tabular}{lp{5.5cm}ll}
    \toprule
    \textbf{Fichero}               & \textbf{Qué se verifica}  & \textbf{VCs} & \textbf{Estado} \\
    \midrule
    \texttt{demo\_verified.java}   & absValue, wrongMethod, sumLoop & 7 & PASS \\
    \texttt{ver01\_simple.java}    & Método sin bucles, @post básico & 2 & PASS \\
    \texttt{ver02\_invariant.java} & Invariante de bucle con variante & 5 & PASS \\
    \bottomrule
  \end{tabular}
\end{table}

\subsection{Salida de ejemplo: \texttt{demo\_verified.java}}

\begin{verbatim}
$ ./compiler --verify tests/verification/demo_verified.java
[OK]  [absValue] correccion (unsat)
[OK]  [sumLoop] iniciacion del invariante (unsat)
[OK]  [sumLoop] preservacion del invariante (unsat)
[OK]  [sumLoop] uso del invariante (unsat)
[OK]  [sumLoop] correccion (unsat)
[OK]  [sumLoop] terminacion (unsat)
[BUG] [wrongMethod] correccion (sat) <-- bug detectado
\end{verbatim}

\section{Script de ejecución}

El script \texttt{tests/run\_tests.sh} automatiza todo el proceso:

\begin{enumerate}
  \item Compila el proyecto con \texttt{make}.
  \item Ejecuta cada prueba de backend: compila el \texttt{.java} con
    \texttt{./compiler}, ejecuta el binario y compara con
    \texttt{expected/}.
  \item Ejecuta las pruebas semánticas comprobando que el compilador
    termina con código $\neq 0$.
  \item Ejecuta las pruebas de verificación comprobando que Z3 esté
    disponible en el \texttt{PATH}.
  \item Imprime un resumen: \texttt{N/M tests passed}.
\end{enumerate}

\begin{verbatim}
$ make && tests/run_tests.sh
[PASS] t01_print
[PASS] t02_arithmetic
[PASS] t03_if_else
[PASS] t04_while
[PASS] t05_functions
[PASS] sem01_undeclared
[PASS] sem02_type_error
[PASS] sem03_redecl
[PASS] ver01_simple
[PASS] ver02_invariant
[PASS] demo_verified
11/11 tests passed
\end{verbatim}


% ── Apéndices ────────────────────────────────────────────────────────────────
\appendix
% ─────────────────────────────────────────────────────────────────────────────
%  Apéndice A – Gramática BNF completa del subconjunto de Java
% ─────────────────────────────────────────────────────────────────────────────
\chapter{Gramática BNF del subconjunto Java}
\label{apend:gramatica}

Las reglas se presentan en notación EBNF. Los nombres entre
\texttt{<>} son no terminales; los literales entre comillas simples
son terminales. Los corchetes denotan opcionalidad; las llaves,
cero o más repeticiones; las alternativas se separan con \texttt{|}.

\section{Estructura del programa}

\begin{lstlisting}[style=pseudo, caption={Producción raíz}]
<Program>  ::= 'public' 'class' ID '{' <Members>* '}'

<Members>  ::= <Annotation>
            |  <Method>

<Annotation> ::= '@pre'       <Expr>
              |  '@post'      <Expr>
              |  '@invariant' <Expr>
              |  '@variant'   <Expr>
\end{lstlisting}

\section{Declaración de métodos}

\begin{lstlisting}[style=pseudo, caption={Métodos y parámetros}]
<Method>  ::= 'public' 'static' <Type> ID '(' [<Params>] ')' <Block>
           |  'public' 'static' 'void' ID '(' [<Params>] ')' <Block>

<Params>  ::= <Type> ID { ',' <Type> ID }

<Type>    ::= 'int' | 'double' | 'boolean' | 'String[]'
\end{lstlisting}

\section{Sentencias}

\begin{lstlisting}[style=pseudo, caption={Sentencias del subconjunto}]
<Block>   ::= '{' <Stmt>* '}'

<Stmt>    ::= <VarDecl>
           |  <Assign>
           |  <If>
           |  <While>
           |  <Return>
           |  <MethodCall> ';'
           |  <PrintStmt>

<VarDecl> ::= <Type> ID ';'
<Assign>  ::= ID '=' <Expr> ';'

<If>      ::= 'if' '(' <Expr> ')' <Block>
           |  'if' '(' <Expr> ')' <Block> 'else' <Block>

<While>   ::= 'while' '(' <Expr> ')' <Annotation>* <Block>

<Return>  ::= 'return' <Expr> ';'

<PrintStmt> ::= 'System.out.println' '(' <Expr> ')' ';'
             |  'System.out.print'   '(' <Expr> ')' ';'
\end{lstlisting}

\section{Expresiones}

\begin{lstlisting}[style=pseudo, caption={Gramática de expresiones (precedencia ascendente)}]
<Expr>  ::= <Expr> '||' <Expr>
         |  <Expr> '&&' <Expr>
         |  '!' <Expr>
         |  <Expr> ('==' | '!=' | '<' | '<=' | '>' | '>=') <Expr>
         |  <Expr> ('+' | '-') <Expr>
         |  <Expr> ('*' | '/' | '%') <Expr>
         |  '-' <Expr>             -- negacion unaria
         |  '(' <Expr> ')'
         |  ID
         |  INT_LIT
         |  REAL_LIT
         |  'true' | 'false'
         |  'result'               -- solo en anotaciones @post
         |  <MethodCall>

<MethodCall> ::= ID '(' [<Args>] ')'
<Args>       ::= <Expr> { ',' <Expr> }
\end{lstlisting}

\section{Anotaciones de verificación}

Las anotaciones de verificación (Milestone~5) son un mecanismo de
extensión léxica: los tokens \texttt{@pre}, \texttt{@post},
\texttt{@invariant} y \texttt{@variant} se reconocen por el analizador
léxico (Flex) y se integran en la gramática como sentencias de primer
nivel dentro de \texttt{Members}, o como prefijos de un \texttt{While}.

Las expresiones de anotación comparten la misma gramática que
\texttt{<Expr>}, con el añadido del símbolo \texttt{result} que
referencia el valor de retorno del método en la postcondición.

\section{Precedencia y asociatividad de operadores}

\begin{table}[H]
  \centering
  \caption{Precedencia (mayor nivel = mayor prioridad)}
  \label{tab:prec}
  \begin{tabular}{cll}
    \toprule
    \textbf{Nivel} & \textbf{Operadores} & \textbf{Asociatividad} \\
    \midrule
    1 & \texttt{||} & izquierda \\
    2 & \texttt{\&\&} & izquierda \\
    3 & \texttt{!} & derecha (prefijo) \\
    4 & \texttt{== !=} & izquierda \\
    5 & \texttt{< <= > >=} & izquierda \\
    6 & \texttt{+ -} & izquierda \\
    7 & \texttt{* / \%} & izquierda \\
    8 & unario \texttt{-} & derecha (prefijo) \\
    9 & llamada a función, literal, identificador & --- \\
    \bottomrule
  \end{tabular}
\end{table}

% ─────────────────────────────────────────────────────────────────────────────
%  Apéndice B – Referencia de instrucciones x86-64
% ─────────────────────────────────────────────────────────────────────────────
\chapter{Referencia de instrucciones x86-64}
\label{apend:instrucciones}

Se documentan las instrucciones x86-64 generadas por el backend del
compilador. Se usa sintaxis AT\&T (fuente a la izquierda, destino a
la derecha).

\section{Instrucciones de transferencia de datos}

\begin{table}[H]
  \centering
  \caption{Instrucciones de movimiento}
  \label{tab:mov}
  \begin{tabular}{lp{8cm}}
    \toprule
    \textbf{Instrucción}        & \textbf{Efecto} \\
    \midrule
    \texttt{movq src, dst}      & $dst \leftarrow src$ (64 bits) \\
    \texttt{movl src, dst}      & $dst \leftarrow src$ (32 bits, zero-extends a 64) \\
    \texttt{movzbq src, dst}    & Extiende byte a 64 bits con ceros (\emph{zero-extend}) \\
    \texttt{leaq src, dst}      & $dst \leftarrow$ dirección efectiva de \textit{src} \\
    \texttt{pushq src}          & Empuja \textit{src} en la pila (\reg{rsp} -= 8) \\
    \texttt{popq dst}           & Extrae de la pila en \textit{dst} (\reg{rsp} += 8) \\
    \bottomrule
  \end{tabular}
\end{table}

\section{Instrucciones aritméticas enteras}

\begin{table}[H]
  \centering
  \caption{Aritmética entera (64 bits)}
  \label{tab:arith-ref}
  \begin{tabular}{lp{8cm}}
    \toprule
    \textbf{Instrucción}        & \textbf{Efecto} \\
    \midrule
    \texttt{addq src, dst}      & $dst \leftarrow dst + src$ \\
    \texttt{subq src, dst}      & $dst \leftarrow dst - src$ \\
    \texttt{imulq src}          & $\reg{rdx}:\reg{rax} \leftarrow \reg{rax} \times src$ (con signo) \\
    \texttt{idivq src}          & $\reg{rax} \leftarrow \reg{rdx}:\reg{rax} \div src$; $\reg{rdx} \leftarrow \text{resto}$ \\
    \texttt{cqto}               & Signo-extiende \reg{rax} a \reg{rdx}:\reg{rax} (necesario antes de \texttt{idivq}) \\
    \texttt{negq dst}           & $dst \leftarrow -dst$ \\
    \texttt{incq dst}           & $dst \leftarrow dst + 1$ \\
    \texttt{decq dst}           & $dst \leftarrow dst - 1$ \\
    \bottomrule
  \end{tabular}
\end{table}

\section{Instrucciones de comparación y salto}

\begin{table}[H]
  \centering
  \caption{Comparación y saltos}
  \label{tab:jmp}
  \begin{tabular}{ll}
    \toprule
    \textbf{Instrucción}        & \textbf{Efecto / Condición} \\
    \midrule
    \texttt{cmpq src, dst}      & Calcula $dst - src$; pone flags (SF, ZF, OF) \\
    \texttt{testq src, dst}     & Calcula $dst \;\mathtt{AND}\; src$; pone ZF \\
    \midrule
    \texttt{je  label}          & Salta si $ZF = 1$ (equal / zero) \\
    \texttt{jne label}          & Salta si $ZF = 0$ \\
    \texttt{jl  label}          & Salta si $SF \neq OF$ (less, con signo) \\
    \texttt{jle label}          & Salta si $ZF = 1$ o $SF \neq OF$ \\
    \texttt{jg  label}          & Salta si $ZF = 0$ y $SF = OF$ (greater) \\
    \texttt{jge label}          & Salta si $SF = OF$ \\
    \texttt{jmp label}          & Salta incondicionalmente \\
    \midrule
    \texttt{sete  al}           & $al \leftarrow (ZF = 1)$ \\
    \texttt{setne al}           & $al \leftarrow (ZF = 0)$ \\
    \texttt{setl  al}           & $al \leftarrow (SF \neq OF)$ \\
    \texttt{setle al}           & $al \leftarrow (ZF \vee SF \neq OF)$ \\
    \texttt{setg  al}           & $al \leftarrow (\neg ZF \wedge SF = OF)$ \\
    \texttt{setge al}           & $al \leftarrow (SF = OF)$ \\
    \bottomrule
  \end{tabular}
\end{table}

\section{Instrucciones de control de flujo y llamadas}

\begin{table}[H]
  \centering
  \caption{Llamadas a función y retorno}
  \label{tab:call}
  \begin{tabular}{lp{8cm}}
    \toprule
    \textbf{Instrucción}          & \textbf{Efecto} \\
    \midrule
    \texttt{call label}           & Empuja \reg{rip}+siguiente y salta a \textit{label} \\
    \texttt{ret}                  & Extrae dirección de retorno de la pila y salta a ella \\
    \texttt{nop}                  & Sin operación (relleno de alineación) \\
    \bottomrule
  \end{tabular}
\end{table}

\section{Instrucciones de punto flotante SSE2}

\begin{table}[H]
  \centering
  \caption{Instrucciones SSE2 para \texttt{double}}
  \label{tab:sse}
  \begin{tabular}{lp{8cm}}
    \toprule
    \textbf{Instrucción}          & \textbf{Efecto} \\
    \midrule
    \texttt{movsd src, dst}       & Mueve un \texttt{double} \\
    \texttt{addsd src, dst}       & $dst \leftarrow dst + src$ (double) \\
    \texttt{subsd src, dst}       & $dst \leftarrow dst - src$ \\
    \texttt{mulsd src, dst}       & $dst \leftarrow dst \times src$ \\
    \texttt{divsd src, dst}       & $dst \leftarrow dst / src$ \\
    \texttt{ucomisd src, dst}     & Comparación de doubles (sin signo, pone ZF/PF/CF) \\
    \texttt{cvtsi2sd src, dst}    & Convierte entero a double \\
    \texttt{cvttsd2si src, dst}   & Convierte double a entero (trunca) \\
    \bottomrule
  \end{tabular}
\end{table}

\section{Directivas del ensamblador}

\begin{table}[H]
  \centering
  \caption{Directivas GAS usadas por el compilador}
  \label{tab:directivas}
  \begin{tabular}{lp{8cm}}
    \toprule
    \textbf{Directiva}            & \textbf{Efecto} \\
    \midrule
    \texttt{.globl sym}           & Exporta el símbolo \textit{sym} (visible al enlazador) \\
    \texttt{.text}                & Sección de código ejecutable \\
    \texttt{.section .rodata}     & Sección de datos de solo lectura (strings de formato) \\
    \texttt{.string "..."}        & Emite una cadena terminada en \verb|\0| \\
    \texttt{.quad n}              & Emite una constante de 8 bytes \\
    \texttt{.align n}             & Alinea la posición actual a $2^n$ bytes \\
    \bottomrule
  \end{tabular}
\end{table}

\section{Convención de registro de marcos (\texttt{\%rbp})}

El compilador usa el marco de pila estándar:

\begin{verbatim}
        Alta dirección
    +------------------+
    | ... args caller  |
    +------------------+ <- %rbp + 16 (primer arg si se pasa en pila)
    | return address   |
    +------------------+ <- %rbp + 8
    | saved %rbp       |
    +------------------+ <- %rbp  <-- base del marco de la funcion
    | local var 0      |
    +------------------+ <- %rbp - 8
    | local var 1      |
    +------------------+ <- %rbp - 16
    | temporales ...   |
    +------------------+ <- %rsp (top de la pila)
        Baja dirección
\end{verbatim}

El tamaño total del marco se calcula como
$\lceil n \cdot 8 / 16 \rceil \times 16$ bytes para mantener
\reg{rsp} alineado a 16 bytes antes de cualquier \texttt{call},
según exige la ABI System V AMD64.


% ── Bibliografía ─────────────────────────────────────────────────────────────
\bibliographystyle{plain}
\bibliography{referencias}

\end{document}
